\documentclass[11pt]{article}
\usepackage{graphicx}
\usepackage{booktabs}
\usepackage{float}
\usepackage{titling}
\setlength{\droptitle}{-3cm}

\title{\textbf{Labwork 6 - Map}}
\author{
Do Minh Quang - 2540095\\
University of Science and Technology of Hanoi
}
\date{}

\begin{document}
\maketitle

\section{GPU Map Implementations}
In this labwork, we will learn how to perform mapping functions inside the CUDA kernel. Three image processing operations are implemented using 2D CUDA grids. The same $718 \times 718$
image as in the previous labwork is used. A second image
of the same size is used for blending.

\begin{figure}[H]
    \centering
    \includegraphics[width=0.48\textwidth]{figure/wc.jpg} \hfill
    \includegraphics[width=0.48\textwidth]{figure/cr.jpg}
    \caption{Input images ($718 \times 718$).}
\end{figure}

For binarization, the RGB values of each pixel are first averaged to obtain a grayscale intensity. A threshold of 128 is then applied: pixels below the threshold are set to 0, while the others are set to 255. The resulting value is written to all three RGB channels.

\begin{figure}[H]
    \centering
    \includegraphics[width=0.4\textwidth]{figure/binary_gpu_2D.jpg}
    \caption{Binary image.}
\end{figure}

For brightness control, a constant value of 100 is added
to or subtracted from each RGB channel to produce brighter
and darker images, respectively. The resulting values
are clipped to the range $[0,255]$ to prevent overflow
and underflow.  

\begin{figure}[H]
    \centering
    \includegraphics[width=0.48\textwidth]{figure/brighter_gpu_2D.jpg} \hfill
    \includegraphics[width=0.48\textwidth]{figure/darker_gpu_2D.jpg}
    \caption{Brightness control images.}
\end{figure}

For image blending, corresponding pixels from two images are combined using
\[
Q(x,y)=c\Phi_1(x,y)+(1-c)\Phi_2(x,y),
\]
where $c=0.5$ in this experiment. 

\begin{figure}[H]
    \centering
    \includegraphics[width=0.4\textwidth]{figure/blend_gpu_2D.jpg}
    \caption{Blending images.}
\end{figure}

\section{Results}
The GPU implementations were measured over 10 runs for each 2D block size. A warm-up run was performed before timing. The measured time includes transferring image data to the GPU and copying the result back to the host.

\begin{table}[H]
\centering
\caption{Average GPU execution time ($\mathrm{ms}\pm\mathrm{std}$) for different 2D block sizes.}
\label{tab:block_sizes}
\begin{tabular}{lccc}
\toprule
Block size & Binarization & Brightness & Blending\\
\midrule
$8\times8$   & $2.2\pm0.6$ & $2.5\pm0.8$ & $3.7\pm0.9$\\
$8\times16$  & $2.3\pm0.6$ & $2.0\pm0.5$ & $3.4\pm0.7$\\
$16\times8$  & $2.3\pm0.6$ & $2.6\pm0.9$ & $3.4\pm0.7$\\
$8\times32$  & $2.6\pm0.9$ & $2.3\pm0.7$ & $3.4\pm0.7$\\
$32\times8$  & $2.3\pm0.6$ & $2.3\pm0.6$ & $3.5\pm0.7$\\
$16\times16$ & $2.3\pm0.6$ & $2.2\pm0.6$ & $3.4\pm0.7$\\
$32\times32$ & $2.4\pm0.6$ & $2.3\pm0.6$ & $3.6\pm0.9$\\
$16\times32$ & $2.3\pm0.6$ & $2.2\pm0.6$ & $3.4\pm0.7$\\
$32\times16$ & $2.3\pm0.6$ & $2.5\pm0.9$ & $3.5\pm0.8$\\
\bottomrule
\end{tabular}
\end{table}

Binarization and brightness control have comparable execution times across the tested block sizes, with their average times generally falling within a similar range at $2.0$--$2.6$ (ms). Blending takes longer, with execution times of $3.4$--$3.7$ (ms), as each thread needs to access two input images instead of one.

The execution time remains relatively stable across
different block sizes, with no consistent improvement
as the block size increases. Binarization is fastest
with $8\times8$ blocks (2.2 ms), while brightness
control is fastest with $8\times16$ blocks (2.0 ms).
For blending, several configurations achieve the
lowest average time of 3.4 ms.

\begin{figure}[H]
\centering
\includegraphics[width=0.78\textwidth]{figure/blocksize_2d_different.png}
\caption{GPU execution time for different 2D block sizes.}
\label{fig:block_size}
\end{figure}

\subsection{Conclusion}
Across the tested configurations, execution times remain relatively stable, with no single block size providing the lowest time for all operations.

\end{document}